\section{Results: IVF Candidate-Generation}
\label{sec:results-ivf}

This section evaluates the approximate two-stage architecture described in
\cref{sec:three-stages}.The architecture consists of IVF token retrieval, 
document selection at budget $C$, and exact float32 MaxSim reranking.The 
results are approximate by construction and are evaluated against the 
exact ranking. We report these findings separately from the exact-safe 
BOND results (\cref{sec:results-bond}).

\subsection{SciFact}
\label{sec:ivf-scifact}

\begin{table*}[t!]
  \centering
  \caption{Controlled SciFact release comparison on 40 held-out queries.
    The configuration uses 5,183 documents, $L=100$, and an $\texttt{nprobe}$ 
    of 8. Values represent medians and the 95th percentile over five interleaved 
    repetitions of the complete online boundary for all 40 queries. Exact R@10 
    the exact ranking recovery against compiled exact MaxSim, and the rerank 
    ratio is the fraction of all query-document pairs evaluated exactly.
    Source: \texttt{scifact\_ivf\_test\_40q\_final.json}.}
  \label{tab:scifact-ivf}
  \footnotesize
  \begin{tabular}{@{}l S[table-format=1.4] S[table-format=1.4] S[table-format=1.4] S[table-format=4.1] S[table-format=1.4]@{}}
    \toprule
    Arm & {Median (\si{\second})} & {p95 (\si{\second})} & {Exact R@10} & {Mean $C$} & {Rerank ratio} \\
    \midrule
    Compiled exact      & 2.5774 & 2.9767 & 1.0000 & 5183.0 & 1.0000 \\
    \midrule
    FAISS-IVF, $C{=}50$   & 0.3550 & 0.3683 & 0.8600 &   50.0 & 0.0096 \\
    PDX-IVF,   $C{=}50$   & 0.3735 & 0.3880 & 0.8600 &   50.0 & 0.0096 \\
    FAISS-IVF, $C{=}100$  & 0.3913 & 0.4183 & 0.9325 &  100.0 & 0.0193 \\
    PDX-IVF,   $C{=}100$  & 0.4257 & 0.4324 & 0.9325 &  100.0 & 0.0193 \\
    FAISS-IVF, $C{=}200$  & 0.4770 & 0.5018 & 0.9750 &  200.0 & 0.0386 \\
    PDX-IVF,   $C{=}200$  & 0.5060 & 0.5149 & 0.9750 &  200.0 & 0.0386 \\
    FAISS-IVF, $C{=}400$  & 0.6949 & 0.7365 & 0.9900 &  390.9 & 0.0754 \\
    PDX-IVF,   $C{=}400$  & 0.7262 & 0.8317 & 0.9900 &  390.9 & 0.0754 \\
    FAISS-IVF, full pool  & 0.9773 & 1.0870 & 1.0000 &  619.3 & 0.1195 \\
    PDX-IVF,   full pool  & 1.0322 & 1.2016 & 1.0000 &  619.3 & 0.1195 \\
    \bottomrule
  \end{tabular}
\end{table*}

Table~\ref{tab:scifact-ivf} shows the held-out SciFact comparison, and 
Figure~\ref{fig:quality-latency} plots the resulting quality--latency 
frontier. Four findings stand out from the data.

\paragraph{Candidate pools are complete; residual loss is selector loss.}
For all 40 held-out queries, the retrieved candidate pool (mean size 619.3 
documents) contains the entire exact top-10. Reranking $C = 200$ documents 
recovers 97.5\% of the exact top-10 while evaluating only 3.86\% of all 
query--document pairs. Reranking the full pool recovers the exact ranking 
with 11.95\% of the pairs.

\paragraph{FAISS-IVF and PDX-IVF are quality-identical.} The artifacts 
show that candidate pools, selected candidates, hit counts, and final top-$k$ 
rankings match between the two engines at every budget. Consequently, their 
recovery curves in Figure~\ref{fig:quality-latency} overlap completely.

\paragraph{No PDX latency advantage.} FAISS has a lower median latency at 
every SciFact budget; the token-search gap ranges from about 0.002 to \SI{0.04}{\second}. 
On NFCorpus, the ordering reverses only at full-pool reranking, where the exact-rerank 
stage happens to be lower for PDX. Across both datasets, the evidence does not support 
a PDX-specific performance benefit.

\paragraph{Qrels metrics cannot resolve the architecture differences.} 
Qrels Recall@10 (0.985) and MRR@10 (0.902) are identical for every arm 
in Table~\ref{tab:scifact-ivf}, even though exact-ranking recovery ranges from 0.86 to 1.00. 
With incomplete relevance judgments and only 40 queries, exact-ranking recovery is the sensitive 
metric; we report qrels metrics only for completeness.

\begin{figure*}[tb]
  \centering
  \includegraphics[width=0.7\textwidth]{fig5_controlled_quality_latency.png}
  \caption{SciFact held-out quality-latency frontier from the controlled
    clean run (5{,}183 documents; 40 queries; $L=100$, $\texttt{nprobe}=8$;
    four pinned physical cores; median of five interleaved repetitions;
    online boundary only). Each point is one rerank budget $C$; the vertical
    axis is recall@10 against the exact MaxSim ranking. FAISS-IVF and PDX-IVF
    trace the same recovery curve, and both reach exact recovery when the full
    candidate pool ($\approx 619$ documents) is reranked, well below the
    compiled exhaustive latency (star).}
  \label{fig:quality-latency}
\end{figure*}

\begin{figure*}[tb]
  \centering
  \includegraphics[width=0.7\textwidth]{fig6_controlled_stage_breakdown.png}
  \caption{Online stage breakdown for the controlled SciFact run (40 held-out
    queries, five interleaved repetitions). Bars stack the median stage timers
    (token search; aggregation and selection; exact rerank and top-$k$)
    against the measured end-to-end median (black marker). At small budgets
    token search dominates; at the full candidate pool exact reranking
    dominates for both engines. Stage medians are diagnostic and need not sum
    exactly to the outer-timer median.}
  \label{fig:stage-breakdown}
\end{figure*}

Figure~\ref{fig:stage-breakdown} decomposes the online time by stage. At 
$C = 50$, token search dominates for both engines; at the full pool, exact 
reranking accounts for most of the latency. Offline index construction, which 
we exclude from these timings, took \SI{18.15}{\second} for FAISS and \SI{21.43}{\second} 
for PDX on SciFact.

\subsection{NFCorpus transfer}
\label{sec:ivf-nfcorpus}

We applied the same stack, online boundary, and configuration 
($L = 100$, $\texttt{nprobe} = 8$, frozen selector) to NFCorpus 
(3,633 documents; 864,703 token vectors; 40 held-out queries). 
Table~\ref{tab:nfcorpus-ivf} reports the results and the left panel of Figure~\ref{fig:transfer} compares the recovery curves of both datasets.

\begin{table*}[tb]
  \centering
  \caption{NFCorpus transfer comparison on 40 held-out queries
    (3{,}633 documents; same configuration and online boundary as
    \cref{tab:scifact-ivf}).
    Source: \texttt{nfcorpus\_ivf\_test\_40q\_final.json}.}
  \label{tab:nfcorpus-ivf}
  \footnotesize
  \begin{tabular}{@{}l S[table-format=1.4] S[table-format=1.4] S[table-format=1.4] S[table-format=4.1] S[table-format=1.4]@{}}
    \toprule
    Arm & {Median (\si{\second})} & {p95 (\si{\second})} & {Exact R@10} & {Mean $C$} & {Rerank ratio} \\
    \midrule
    Compiled exact      & 1.2554 & 1.4940 & 1.0000 & 3633.0 & 1.0000 \\
    \midrule
    FAISS-IVF, $C{=}50$  & 0.2268 & 0.2585 & 0.8275 &   50.0 & 0.0138 \\
    PDX-IVF,   $C{=}50$  & 0.2334 & 0.2648 & 0.8275 &   50.0 & 0.0138 \\
    FAISS-IVF, $C{=}200$ & 0.3460 & 0.3539 & 0.9650 &  198.0 & 0.0545 \\
    PDX-IVF,   $C{=}200$ & 0.3598 & 0.4300 & 0.9650 &  198.0 & 0.0545 \\
    FAISS-IVF, full pool & 0.5698 & 0.5771 & 0.9750 &  388.1 & 0.1068 \\
    PDX-IVF,   full pool & 0.5358 & 0.6432 & 0.9750 &  388.1 & 0.1068 \\
    \bottomrule
  \end{tabular}
\end{table*}

The transfer is less forgiving than SciFact in one respect: reranking the 
entire candidate pool reaches only 0.975 exact recall@10, because the pool 
contains the complete exact top-10 for only 37 of the 40 queries. The residual 
0.025 loss is therefore candidate-pool coverage loss, meaning that no selector 
operating on this pool can repair it. On SciFact, by contrast, all loss at 
finite $C$ was selector loss. Candidate pools and final rankings again agree 
between FAISS and PDX at every budget.

The PDX token-search stage is slightly slower than FAISS in every matched 
NFCorpus arm (for example, median \SI{0.155}{\second} versus \SI{0.142}{\second} at 
$C = 50$), but total online latency is nearly tied at $C = 200$ and at the 
full pool. Thus the earlier conclusion holds: the evidence shows no consistent 
PDX-specific advantage. Qrels Recall@10 moves only from 0.2150 at $C=50$ to 
0.2191 at full pool, which equals the exact reference, confirming again that 
qrels cannot resolve these differences. Offline index construction took 
\SI{12.37}{\second} for FAISS and \SI{14.25}{\second} for PDX.